\begin{lemma}
\label{editorial-source-lemma-characterize-jacobson}
Let $R$ be a ring. If $R$ is not Jacobson there exist
a prime $\mathfrak p \subset R$, an element $f \in R$
such that the following hold
\begin{enumerate}
\item $\mathfrak p$ is not a maximal ideal,
\item $f \not \in \mathfrak p$,
\item $V(\mathfrak p) \cap D(f) = \{\mathfrak p\}$, and
\item $(R/\mathfrak p)_f$ is a field.
\end{enumerate}
On the other hand, if $R$ is Jacobson, then for any pair $(\mathfrak p, f)$
such that (1) and (2) hold, $V(\mathfrak p)\cap D(f)$ contains
infinitely many maximal ideals of $R$; in particular it is infinite.
\end{lemma}

\begin{proof}
Assume $R$ is not Jacobson.
By Lemma \ref{algebra-lemma-jacobson} this means there exists a
closed subset $T \subset \Spec(R)$
whose set $T_0 \subset T$ of closed points is not dense in $T$.
Choose an $f \in R$ such that $T_0 \subset V(f)$ but
$T \not \subset V(f)$. Note that $T \cap D(f)$
is homeomorphic to $\Spec((R/I)_f)$ if $T = V(I)$, see
Lemmas \ref{algebra-lemma-spec-closed} and \ref{algebra-lemma-standard-open}.
The set $T\cap D(f)$ is nonempty, so its displayed coordinate ring
is nonzero. Every nonzero ring has a maximal ideal
(Lemma \ref{editorial-source-lemma-Zariski-topology}), so choose a closed point $t$
of the space $T\cap D(f)$. Then $t$ corresponds to a prime ideal
$\mathfrak p \subset R$ which is not maximal (as $t \not \in T_0$).
Thus (1) holds. By construction $f \not \in \mathfrak p$, hence (2).
As $t$ is a closed point of $T \cap D(f)$ we see that
$V(\mathfrak p) \cap D(f) = \{\mathfrak p\}$, i.e., (3) holds. Hence we
conclude that $(R/\mathfrak p)_f$ is a domain whose
spectrum has one point, hence (4) holds
(for example combine Lemmas \ref{editorial-source-lemma-characterize-local-ring} and
\ref{algebra-lemma-minimal-prime-reduced-ring}).

\medskip\noindent
Conversely, suppose $R$ is Jacobson and $(\mathfrak p,f)$ satisfies
(1) and (2). Suppose only finitely many maximal ideals
$\mathfrak m_1,\ldots,\mathfrak m_t$ lie in
$V(\mathfrak p)\cap D(f)$. Because $\mathfrak p$ is not maximal,
choose $g_i\in\mathfrak m_i\setminus\mathfrak p$ and let
$g=\prod_{i=1}^t g_i$, with $g=1$ when $t=0$.
Primality gives $fg\notin\mathfrak p$, so the locally closed subset
$V(\mathfrak p)\cap D(fg)$ is nonempty. Since $\Spec(R)$ is Jacobson,
it contains a point closed in $\Spec(R)$. That point would be one of
the listed maximal ideals, since $D(fg)\subset D(f)$, but $g$ belongs
to every listed ideal, excluding each of them. This contradiction
proves that there are infinitely many such maximal ideals.
\end{proof}